\documentclass[11pt]{article}
\usepackage{amsmath, amssymb}
\usepackage[margin=1in]{geometry}
\usepackage{fancyhdr}

% ---- Fill these in ----
\newcommand{\courseName}{NEU 314}
\newcommand{\taName}{Keith Murray}
\newcommand{\sectionInfo}{Precept B01}
% ------------------------

\title{Precept 2: Warm Up Exercise}
\date{\today}

\pagestyle{fancy}
\fancyhf{}
\lhead{\courseName}
\chead{\sectionInfo}
\rhead{\today}
\cfoot{\thepage}

\begin{document}

\maketitle
\thispagestyle{fancy}

\section*{Problem}

Consider the following vectors:
\[
\vec{v}=\begin{bmatrix}
    2\\2\\0
\end{bmatrix},\qquad\vec{u}_1=\begin{bmatrix}
    1\\2\\3
\end{bmatrix},\qquad\vec{u}_2=\begin{bmatrix}
    1\\1\\-1
\end{bmatrix}
\]
Let's answer the following questions:
\begin{enumerate}
    \item What is the relationship between $\vec{u}_1$ and $\vec{u}_2$?
    \item What is the dimension of the subspace spanned by $\vec{u}_1$ and $\vec{u}_2$?
    \item How would you project vector $\vec{v}$ on to the subspace spanned by $\vec{u}_1$ and $\vec{u}_2$?
\end{enumerate}
Remember the formula for vector projection is $\operatorname{proj}_{\vec{u}}\left(\vec{v}\right)=\frac{\vec{v}\cdot \vec{u}}{\vec{u}\cdot \vec{u}}\vec{u}$

\section*{Solution}
\subsection*{Question 1}
$\vec{u}_1$ and $\vec{u}_2$ are orthogonal because
\[
\vec{u}_1\cdot \vec{u}_2 = 1\cdot 1 + 2\cdot 1 + 3\cdot\left(-1\right)=3-3=0
\]

\subsection*{Question 2}
Note that nonzero orthogonal vectors are automatically linearly independent. Hence, the subspace is a 2-dimensional subspace because there are two vectors and they are linearly independent.

\subsection*{Question 3}

Because $\vec{u}_1$ and $\vec{u}_2$ are orthogonal, we can use the vector projection formula on both $\vec{v}$ on to $\vec{u}_1$ and $\vec{v}$ on to $\vec{u}_2$, and then add the resulting projections. Mathematically, this looks like
\[
\frac{\vec{v}\cdot \vec{u}_1}{\vec{u}_1\cdot \vec{u}_1}\vec{u}_1 + \frac{\vec{v}\cdot \vec{u}_2}{\vec{u}_2\cdot \vec{u}_2}\vec{u}_2
\]
which evaluates to 
\[
\frac{3}{7}\vec{u}_1 + \frac{4}{3}\vec{u}_2=\frac{1}{21}\begin{bmatrix}
    37\\46\\-1
\end{bmatrix}
\]
Note that this would not work if $\vec{u}_1$ and $\vec{u}_2$ were not orthogonal.


\end{document}